\subsection{Bayseian Inversion}

% Breif introduction

By applying the inversion framework outlined in \autoref{sec:bay} we can combined our understanding of SL physics and with prior distributions built to reflect our knowledge of the system to result in an updated posterior belief of the state of the model space based on observations made in the data space.



\subsubsection{Simple inversion}

\notebookcite{thomasholland06SimpleInversionipynb2026}

A simple case for this is to invert for ice thickness change, $\thicknesschange{i}$, with a single density, from SSH altimetery, $d$ where the forward problem takes the form

\begin{align}
	d &= \operatoraltimetry{\text{ocean}} \circ \operatorsshc \circ \operatorslc \circ \operatorloadice (\thicknesschange{i}) + n,\\
	d &= \operator{O}{simple}(\thicknesschange{i}) + n, \label{eq:simple_case}
\end{align}

where $\operator{O}{simple}$ is the composed forward operator for this simple case and $n$ represents the altimetry noise. The noise is taken to be an independent Gaussian measure in the dataspace with zero mean and covariance $\covariance{n}$, so

\begin{equation}
	n \sim \mathcal{N}(0, \covariance{n}).
\end{equation}

A true ice-thickness field, $\thicknesschange{i}^\dagger$,  is drawn as a single sample from the prior measure $\thicknesschange{i}$, providing a plausible but non-trivial test case that is, by construction, consistent with our prior assumptions. Synthetic altimetry observations, $d_{obs}$, are then generated by passing through the forward operator and adding the Gaussian noise:

\begin{equation}
	d_{obs} = \operator{O}{simple}(\thicknesschange{i}^\dagger) + n
\end{equation}

Using the Bayesian inversion framework, we are able to generate the posterior model space measure, $m_{\Delta I}^{p}$, which can be inspected to understand the new inversion method. Such inspections can include simple visual inspection of the posterior model expectation compared to the true field is is inverting for. Both of these can be pushed forward through the sea level operator to obtain their respective sea level fingerprints for comparison.

The covariance between the major sources can be inspected to understand how well the method is able to disambiguate between them.

\subsubsection{Comparison with previous method}
\notebookcite{thomasholland06SimpleInversionipynb2026}

To compare our method's estimate of GMSLC with the altimetry estimate, we can compare the $z$-scores of the two methods' estimates of GMSL change relative to the true GMSL change. We generate a true GMSL value, $\gmslctrue$, and the inversion method's estimate, $\gmslcbay$ by applying the mass balance GMSL operator, $\operatorgmsl{\text{true}}$, to the true ice thickness change field and the posterior model expectation, respectively:

\begin{align}
	\gmslctrue &= \operatorgmsl{\text{true}} \thicknesschange{i}^\dagger, \label{eq:true-gmsl-model} \\
	\gmslcbay &= \operatorgmsl{\text{true}} \thicknesschange{i}^{p}. \label{eq:bay-gmsl-model}
\end{align}


The GMSL altimetry estimate is calculated from the available sea surface height change (SSHC) point measurements, $d_i$, using an area-weighted scheme (\autoref{eq:gmsl_ssh_alt_integral}). Each individual observation is treated as an independent 1D Gaussian measure

\begin{equation}
  d_i \sim \mathcal{N}(d_i, \covariance{n,\,i}),
\end{equation}

where $\covariance{n,\,i}$ represents the noise variance at that specific point. These individual measures are then combined via a cosine-latitude weighted average (with weights $w_i$) to form a single global distribution. Mathematically, this estimate represents the true ice state evaluated through the forward operators plus an aggregated noise term, yielding the following distribution:

\begin{align}
 \gmslcalt &= \operatorgmsl{\text{point}} \circ \operatorsshc \circ \operatorslc \circ \operatorloadice (\thicknesschange{i}^\dagger) + n_{\text{alt}} \\
 &\sim \mathcal{N}\left( \sum_{i=1}^{N} w_i d_i, \sum_{i=1}^{N} w_i^2 \covariance{n,\,i} \right) \label{eq:alt-gmsl-model}
\end{align}

The $z$-score is calculated using the expectation and standard deviations, $\mu$ and $\sigma$ respectively, for both the old altimetry method and the new inversion method, to compare the performance of the two methods in estimating the true GMSL change:

\begin{align}
    z^{bay} &= \frac{\gmslctrue - \mu_{\text{bay}}}{\sigma_{\text{bay}}} \\[7pt]
    z^{alt} &= \frac{\gmslctrue - \mu_{\text{alt}}}{\sigma_{\text{alt}}}
\end{align}

The inversions was run multiple times to generate a suite of $z$-scores for each method, which can be compared to show the improvement of the new method over the old method.

\subsubsection{Sensitivity tests}
\notebookcite{thomasholland07InversionSensitivityipynb2026}

With any inverse setup, it is important to understand how sensitive the results are to the various parameters of the method. This is particularly significant for our method, as it relies on a number of assumptions and choices in the construction of the prior model space measure.

To test the inversion method, the scheme described in \textcite{heathcoteScaleableBayesianFramework2026} for testing the inversion sensitivity to changes in the ice lengthscale, covariance, and expectation bias parameters of the prior.

Initially, an inversion scenario is set up similar to the inversion, with a fixed truth ice state, $m_{\Delta I}^\dagger$. Inversions for the true field were run with priors generated with the systematically varied parameters of interest, and comparison is made between their results and the true field to assess the method's sensitivity to that parameter.

In addition to the above tests from \textcite{heathcoteScaleableBayesianFramework2026}, the sensitivity of the method to the ice thickness prior model to spatial patters or not. True ice states for both spatially uniform ice change and thickness dependent change patters were generated (see \autoref{sec:ice_thickness_priors}). These were both inverted for using the two types of prior model, and then compared.

\subsubsection{Joint inversion using ice, firn, and ocean dynamic changes}
% Paragraph: Describe the joint inversion -- ice + firn + ocean dynamics
% inverted simultaneously. Explain why this matters (correlated signals).

The model space would be a direct sum of the ice thickness change, firn thickness change space, and ocean dynamic topography change:

\begin{equation}
	m = \thicknesschange{i} \oplus \thicknesschange{f} \oplus \thicknesschange{odt}
\end{equation}

With the data space being the direct sum of SSH altimetry data, tide gauge data, and ice altimetry data:

\begin{equation}
	d = d_{\Delta SSH} \oplus d_{\Delta TG} \oplus d_{\Delta ISH}
\end{equation}

It is often the case where multiple model space parameters map to one data space parameter, for example in our case, SSH altimetry points forms the linear equation:

\begin{equation}
	d_{\Delta SSH} = \operatorpointocean( \operatorsshc \circ \operatorslc (\operatorloadice m_{\Delta I} + \operatorloadfirn m_{\Delta F} + \operatorloadwater \odtchange) + \odtchange)
\end{equation}

The forward operator for the joint inversion can be written as a block matrix, with each block being the mapping of the respective model space parameter to the respective data space parameter. This results in the following form:

\begin{align}
\begin{bmatrix}
	\text{SSH Altimetry Points}\\
	\text{Tide Gauge Points}\\
	\text{Ice Altimetry Points}
\end{bmatrix}
&\nonumber\\&=
\begin{bmatrix}
	\operatorpointocean \operatorsshc \operatorslc \operatorloadice & \operatorpointocean \operatorsshc \operatorslc \operatorloadfirn & \operatorpointocean(\operatorsshc \operatorslc \operatorloadwater + 1) \\
	\operatorpointtg \operatorslc \operatorloadice & \operatorpointtg \operatorslc \operatorloadice & \operatorpointtg(\operatorslc \operatorloadwater + 1)\\
	\operatorpointice & \operatorpointice & \mathbf{0}
\end{bmatrix}
\begin{bmatrix}
	\thicknesschange{\text{Ice}}\\
	\thicknesschange{\text{Firn}}\\
	\odtchange
\end{bmatrix}
\end{align}

However, this is computationally inefficient, as it requires six independent calls to the fingerprint numerical solver. To optimise this forward operator to reduce the number of fingerprint calls, we can factorise this matrix into sparse matrices (see \autoref{sec:joint_operator_factorisation}), shown in \autoref{eq:factorised_joint}. This reduced each iteration time, with a speed-up of $3\times$.

\begin{equation}
=
	\begin{bmatrix}
\operatorpointocean \operatorsshc & \operatorpointocean & 0 & 0 \\
\operatorpointtg & \operatorpointtg & 0 & 0 \\
0 & 0 & \operatorpointice & \operatorpointice
\end{bmatrix}
\begin{bmatrix}
\operatorslc & 0 & 0 & 0 \\
0 & \mathbf{1} & 0 & 0 \\
0 & 0 & \mathbf{1} & 0 \\
0 & 0 & 0 & \mathbf{1}
\end{bmatrix}
\begin{bmatrix}
\mathbf{L_I} & \mathbf{L_F} & \mathbf{L_W} \\
0 & 0 & \mathbf{1} \\
\mathbf{1} & 0 & 0 \\
0 & \mathbf{1} & 0
\end{bmatrix}
\begin{bmatrix}
\thicknesschange{\text{Ice}} \\
\thicknesschange{\text{Firn}} \\
\odtchange
\end{bmatrix}\label{eq:factorised_joint}
\end{equation}

% Paragraph: Discuss preconditioning and why it is needed for numerical
% efficiency when the operator is ill-conditioned.


% Paragraph: Describe application to altimetry data and/or DUACS SSH data.

\subsubsection{Sensitivity tests of joint inversions}

TKTK covariance tests

TKTK knock out tests
